\documentclass[../../main.tex]{subfiles}
\begin{document}

\section{光滑结构}\label{section:ch01-光滑结构}

我们在上一节给出的流形定义足以研究流形的拓扑性质，如紧性、连通性、单连通性以及流形在同胚下的分类问题。然而，在整个拓扑流形理论中完全没有提到微积分。这有一个充分的理由：无论我们怎样设法使流形上函数的导数有意义，这类导数都不可能在同胚下不变。例如，由
\begin{align*}
\varphi(u,v)=(u^{1/3},v^{1/3})
\end{align*}
给出的映射 $\varphi:\mathbb{R}^{2}\to\mathbb{R}^{2}$ 是一个同胚，而且容易构造可微函数 $f:\mathbb{R}^{2}\to\mathbb{R}$，使得 $f\circ\varphi$ 在原点不可微。（函数 $f(x,y)=x$ 就是这样一个例子。）

为了使实值函数、曲线或流形之间映射的导数有意义，我们需要引入一种新的流形，称为光滑流形。它将是一个拓扑流形，除拓扑外还带有某种额外结构，使我们能够判定哪些从流形出发或映到流形的函数是光滑的。其定义将建立在欧几里得空间之间映射的微积分之上，因此我们先回顾一下关于这类映射的一些基本术语。若 $U$ 和 $V$ 分别是欧几里得空间 $\mathbb{R}^{n}$ 和 $\mathbb{R}^{m}$ 的开子集，则称函数 $F:U\to V$ 是光滑的（或 $C^{\infty}$ 的，或无穷次可微的），如果它的每个分量函数都具有各阶连续偏导数。若此外 $F$ 是双射且具有光滑逆映射，则称它为微分同胚。微分同胚特别地是一个同胚。

光滑映射的一些重要性质在附录~C 中给出回顾。你应当注意，有些作者对"光滑"一词的定义不同——例如，取为连续可微或仅仅可微。另一方面，有些作者用"可微"一词表示我们所说的光滑。在本书中，光滑与 $C^{\infty}$ 同义。

为了看出拓扑流形上什么样的额外结构适合用来辨别哪些映射是光滑的，考虑任意一个拓扑 $n$ 流形 $M$。$M$ 中的每个点都在某个坐标映射 $\varphi:U\to\widetilde{U}\subset\mathbb{R}^{n}$ 的定义域中。关于 $M$ 上光滑函数的一个合理定义是说：$f:M\to\mathbb{R}$ 光滑当且仅当复合函数 $f\circ\varphi^{-1}:\widetilde{U}\to\mathbb{R}$ 在通常微积分意义下光滑。但这只有在该性质与坐标图的选取无关时才有意义。为保证这种无关性，我们将把注意力限制在"光滑图"上。由于光滑性不是同胚不变性质，做到这一点的方法是把所有光滑图的集合看作 $M$ 上的一种新结构。

带着这一动机，我们现在来描述构造的细节。设 $M$ 是一个拓扑 $n$ 流形。若 $(U,\varphi)$、$(V,\psi)$ 是两个满足 $U\cap V\neq\varnothing$ 的图，则复合映射 $\psi\circ\varphi^{-1}:\varphi(U\cap V)\to\psi(U\cap V)$ 称为从 $\varphi$ 到 $\psi$ 的转移映射（图~\cref{figure:ch01-1.6}）。它是同胚的复合，因而本身也是一个同胚。称两个图 $(U,\varphi)$ 与 $(V,\psi)$ 是光滑相容的，如果 $U\cap V=\varnothing$，或者转移映射 $\psi\circ\varphi^{-1}$ 是一个微分同胚。由于 $\varphi(U\cap V)$ 与 $\psi(U\cap V)$ 是 $\mathbb{R}^{n}$ 的开子集，该映射的光滑性应按通常意义理解，即具有各阶连续偏导数。

\begin{figure}[ht]
\centering
\includegraphics[width=0.8\textwidth]{fig-01-6.png}
\caption{转移映射}\label{figure:ch01-1.6}
\end{figure}

我们定义 $M$ 的一个图册为定义域覆盖 $M$ 的一族图。称一个图册 $\mathcal{A}$ 为光滑图册，如果 $\mathcal{A}$ 中任意两个图彼此光滑相容。

为证明一个图册是光滑的，只需验证每当 $(U,\varphi)$ 和 $(V,\psi)$ 是 $\mathcal{A}$ 中的图时转移映射 $\psi\circ\varphi^{-1}$ 光滑；一旦证明了这一点，便可知 $\psi\circ\varphi^{-1}$ 是微分同胚，因为其逆 $(\psi\circ\varphi^{-1})^{-1}=\varphi\circ\psi^{-1}$ 正是我们已经证明光滑的转移映射之一。或者，给定两个具体的图 $(U,\varphi)$ 和 $(V,\psi)$，通常最容易的做法是验证 $\psi\circ\varphi^{-1}$ 光滑、单射且在每点处 Jacobian 非奇异，然后引用推论~C.36。

我们的计划是通过给出一个光滑图册来定义 $M$ 上的"光滑结构"，并定义函数 $f:M\to\mathbb{R}$ 光滑当且仅当对图册中每个坐标图 $(U,\varphi)$，$f\circ\varphi^{-1}$ 在通常微积分意义下光滑。这个做法有一个小小的技术问题：一般而言，会有许多可能的图册给出"同一个"光滑结构，即它们都确定 $M$ 上相同的光滑函数集合。例如，考虑 $\mathbb{R}^{n}$ 上如下一对图册：
\begin{align*}
\mathcal{A}_{1}&=\{(\mathbb{R}^{n},\operatorname{Id}_{\mathbb{R}^{n}})\},\\
\mathcal{A}_{2}&=\{(B_{1}(x),\operatorname{Id}_{B_{1}(x)}):x\in\mathbb{R}^{n}\}.
\end{align*}
尽管它们是不同的光滑图册，但显然函数 $f:\mathbb{R}^{n}\to\mathbb{R}$ 关于其中任意图册光滑，当且仅当它在通常微积分意义下光滑。

我们可以选择把光滑结构定义为在适当等价关系下光滑图册的等价类。然而，更直接的是给出如下定义：称 $M$ 上的一个光滑图册 $\mathcal{A}$ 是极大的，如果它不被任何更大的光滑图册真包含。这恰好意味着每个与 $\mathcal{A}$ 中每个图都光滑相容的图已经在 $\mathcal{A}$ 中。（这样的光滑图册也称为完全的。）

现在我们可以定义本章的主要概念。若 $M$ 是一个拓扑流形，则 $M$ 上的一个光滑结构是一个极大光滑图册。光滑流形是一对 $(M,\mathcal{A})$，其中 $M$ 是一个拓扑流形，$\mathcal{A}$ 是 $M$ 上的一个光滑结构。当光滑结构不言自明时，我们通常略去对它的提及，而只说"$M$ 是一个光滑流形"。光滑结构也被某些作者称为可微结构或 $C^{\infty}$ 结构。我们还用术语光滑流形结构来指一个流形拓扑连同一个光滑结构。

我们强调，光滑结构是必须加到拓扑流形上的一份额外数据，此后我们才有权谈论"光滑流形"。事实上，一个给定的拓扑流形可能有许多不同的光滑结构（见例~\cref{example:ch01-另一光滑结构} 和习题~1-6）。另一方面，并非总能在给定的拓扑流形上找到光滑结构：存在根本不允许任何光滑结构的拓扑流形。（第一个例子是 Michel Kervaire 于 1960 年发现的一个 $10$ 维紧流形 [Ker60]。）

通过显式描述一个极大光滑图册来定义光滑结构通常不太方便，因为这样的图册包含非常多的图。幸运的是，如下一个命题所示，我们只需指定某个光滑图册即可。

\begin{proposition}\label{proposition:ch01-极大图册唯一}
设 $M$ 是一个拓扑流形。
\begin{enumerate}[(a)]
\item $M$ 的每个光滑图册 $\mathcal{A}$ 都包含于唯一的一个极大光滑图册中，该极大光滑图册称为由 $\mathcal{A}$ 确定的光滑结构。
\item $M$ 的两个光滑图册确定同一个光滑结构，当且仅当它们的并是一个光滑图册。
\end{enumerate}
\end{proposition}

\begin{proof}
设 $\mathcal{A}$ 是 $M$ 的一个光滑图册，令 $\mathcal{S}_{\mathcal{A}}$ 为所有与 $\mathcal{A}$ 中每个图都光滑相容的图所成的集合。为证明 $\mathcal{S}_{\mathcal{A}}$ 是一个光滑图册，我们需要证明 $\mathcal{S}_{\mathcal{A}}$ 中任意两个图彼此光滑相容，也就是说对任意 $(U,\varphi),(V,\psi)\in\mathcal{S}_{\mathcal{A}}$，映射 $\psi\circ\varphi^{-1}:\varphi(U\cap V)\to\psi(U\cap V)$ 光滑。

任取 $x=\varphi(p)\in\varphi(U\cap V)$。因为 $\mathcal{A}$ 中图的定义域覆盖 $M$，存在某个图 $(W,\theta)\in\mathcal{A}$ 使 $p\in W$（图~\cref{figure:ch01-1.7}）。由于 $\mathcal{S}_{\mathcal{A}}$ 中每个图都与 $(W,\theta)$ 光滑相容，映射 $\theta\circ\varphi^{-1}$ 和 $\psi\circ\theta^{-1}$ 在其定义处都光滑。因为 $p\in U\cap V\cap W$，由此可得
\begin{align*}
\psi\circ\varphi^{-1}=(\psi\circ\theta^{-1})\circ(\theta\circ\varphi^{-1})
\end{align*}
在 $x$ 的一个邻域上光滑。因此，$\psi\circ\varphi^{-1}$ 在 $\varphi(U\cap V)$ 中每一点的邻域上都光滑。所以 $\mathcal{S}_{\mathcal{A}}$ 是一个光滑图册。为验证它是极大的，只需注意任何与 $\mathcal{S}_{\mathcal{A}}$ 中每个图都光滑相容的图，特别地与 $\mathcal{A}$ 中每个图光滑相容，因而已经在 $\mathcal{S}_{\mathcal{A}}$ 中。这证明了包含 $\mathcal{A}$ 的极大光滑图册的存在性。若 $\mathcal{B}$ 是包含 $\mathcal{A}$ 的任意另一个极大光滑图册，则它的每个图都与 $\mathcal{A}$ 中每个图光滑相容，故 $\mathcal{B}\subset\mathcal{S}_{\mathcal{A}}$。由 $\mathcal{B}$ 的极大性，$\mathcal{B}=\mathcal{S}_{\mathcal{A}}$。

(b) 的证明留作练习。
\end{proof}

\begin{figure}[ht]
\centering
\includegraphics[width=0.8\textwidth]{fig-01-7.png}
\caption{命题 1.17(a) 的证明}\label{figure:ch01-1.7}
\end{figure}

\begin{exercise}\label{exercise:ch01-1.18}
证明命题~\cref{proposition:ch01-极大图册唯一}(b)。
\end{exercise}

例如，若拓扑流形 $M$ 可由单个坐标图覆盖，则光滑相容条件平凡地满足，所以任意这样的图都自动在 $M$ 上确定一个光滑结构。

值得一提的是，光滑结构的概念可以通过改变对图的相容要求而以几种不同方式推广。例如，如果我们把图光滑相容的要求替换为较弱的要求，即每个转移映射 $\psi\circ\varphi^{-1}$（及其逆）属于 $C^{k}$ 类，便得到 $C^{k}$ 结构的定义。类似地，如果我们要求每个转移映射是实解析的（即可在每点的邻域内表示为收敛幂级数），便得到实解析结构的定义，也称为 $C^{\omega}$ 结构。若 $M$ 有偶维数 $n=2m$，我们可以把 $\mathbb{R}^{2m}$ 等同于 $\mathbb{C}^{m}$，并要求转移映射是复解析的；这确定了一个复解析结构。赋予其中一种结构的流形分别称为 $C^{k}$ 流形、实解析流形或复流形。（注意 $C^{0}$ 流形恰好就是拓扑流形。）本书不讨论这些其他种类的流形，但它们在分析中起重要作用，因此了解这些定义是有用的。

\subsection{局部坐标表示}\label{subsection:ch01-局部坐标表示}

若 $M$ 是一个光滑流形，则包含在给定极大光滑图册中的任意图 $(U,\varphi)$ 称为光滑图，相应的坐标映射 $\varphi$ 称为光滑坐标映射。引入术语光滑坐标域或光滑坐标邻域来指光滑坐标图的定义域也是有用的。光滑坐标球指其在光滑坐标映射下的像为欧几里得空间中一个球的光滑坐标域。光滑坐标立方体类似地定义。

通常有用的是把注意力限制在那些闭包漂亮地落在更大坐标球内部的坐标球上。称集合 $B\subset M$ 为正则坐标球，如果存在一个光滑坐标球 $B'\supset\overline{B}$ 和一个光滑坐标映射 $\varphi:B'\to\mathbb{R}^{n}$，使得对某对正实数 $r<r'$ 有
\begin{align*}
\varphi(B)=B_{r}(0),\qquad \varphi(\overline{B})=\overline{B_{r}(0)},\qquad \varphi(B')=B_{r'}(0).
\end{align*}
因为 $\overline{B}$ 同胚于 $\overline{B_{r}(0)}$，它是紧的，因此每个正则坐标球在 $M$ 中都是预紧的。下一个命题对光滑流形给出了引理~\cref{lemma:ch01-预紧坐标球可数基} 的一个稍许改进。其证明是该引理证明的直接改写。

\begin{proposition}\label{proposition:ch01-正则坐标球可数基}
每个光滑流形都有一个由正则坐标球组成的可数基。
\end{proposition}

\begin{exercise}\label{exercise:ch01-1.20}
证明命题~\cref{proposition:ch01-正则坐标球可数基}。
\end{exercise}

下面是人们通常对光滑流形上坐标图的思考方式。一旦我们在 $M$ 上选取了一个光滑图 $(U,\varphi)$，坐标映射 $\varphi:U\to\widetilde{U}\subset\mathbb{R}^{n}$ 就可以看作给出了 $U$ 与 $\widetilde{U}$ 之间的一个临时等同。利用这个等同，当我们在这个图中工作时，可以同时把 $U$ 看作 $M$ 的一个开子集和 $\mathbb{R}^{n}$ 的一个开子集。你可以通过想象画在 $U$ 上的一张"网格"来直观理解这个等同，该网格代表坐标直线在 $\varphi$ 下的原像（图~\cref{figure:ch01-1.8}）。在这个等同下，我们可以用 $p$ 的坐标 $(x^{1},\dots,x^{n})=\varphi(p)$ 来表示 $p\in U$，并把这个 $n$ 元组看作点 $p$。我们通常这样表达："$(x^{1},\dots,x^{n})$ 是 $p$ 的（局部）坐标表示"，或"$p=(x^{1},\dots,x^{n})$ 在局部坐标下"。

\begin{figure}[ht]
\centering
\includegraphics[width=0.8\textwidth]{fig-01-8.png}
\caption{坐标网格}\label{figure:ch01-1.8}
\end{figure}

另一种看法是，借助我们的等同 $U\leftrightarrow\widetilde{U}$，可以把 $\varphi$ 看作恒等映射并从记号中略去它。这需要一点适应，但回报是在许多情形下记号的极大简化。你只需记住这个等同一般只是局部的，并且强烈依赖于坐标图的选取。

你大概在多元微积分的学习中已经习惯了这样的等同。最常见的例子是平面中的极坐标 $(r,\theta)$，由关系 $(x,y)=(r\cos\theta,r\sin\theta)$ 隐式定义（见例~C.37）。在一个适当的开子集上，例如 $U=\{(x,y):x>0\}\subset\mathbb{R}^{2}$，$(r,\theta)$ 可以表示为 $(x,y)$ 的光滑函数，而把 $(x,y)$ 映到相应 $(r,\theta)$ 的映射是关于 $\mathbb{R}^{2}$ 上标准光滑结构的一个光滑坐标映射。利用这个映射，给定的点 $p\in U$ 既可以在标准坐标下写为 $p=(x,y)$，也可以在极坐标下写为 $p=(r,\theta)$，其中两个坐标表示由
\begin{align*}
(r,\theta)=\bigl(\sqrt{x^{2}+y^{2}},\tan^{-1}(y/x)\bigr),\qquad (x,y)=(r\cos\theta,r\sin\theta)
\end{align*}
相联系。其他极坐标图可以通过把 $(r,\theta)$ 限制到 $\mathbb{R}^{2}\setminus\{0\}$ 的其他开子集来得到。

流形不附带任何预先指定的坐标选择这一事实，既是福也是祸。或多或少任意地选择坐标的灵活性，在处理流形理论问题时可能是一大优势，因为坐标往往可以选得使手头问题的某个方面得到简化。但我们为这种灵活性付出的代价是：必须保证我们希望在流形上整体定义的任何对象都不依赖于坐标的特定选取。做到这一点通常有两种方式：要么写下一个依赖坐标的定义，然后证明该定义在任意坐标图中都给出相同结果；要么写下一个显然与坐标无关的定义（常称为不变定义）。在少数情形下依赖坐标的方法明显更简单，我们会采用，但在大多数情况下我们会尽可能给出与坐标无关的定义。对这类定义的需求，是现代流形理论抽象性的主要来源之一。为了有效使用流形理论，你需要培养的最重要技能之一，就是能够在不变描述及其坐标对应物之间轻松来回切换。

\end{document}